% Archivo generado por bd/generar_tablas.ps1 a partir de bd/fase_posterior/crear_tablas_fase_posterior.sql. No editar a mano.
\begin{center}\small
\begin{longtable}{|L{0.19\textwidth}|L{0.23\textwidth}|L{0.08\textwidth}|L{0.38\textwidth}|}
\hline
\textbf{Entidad} & \textbf{Llaves candidatas} & \textbf{Forma} & \textbf{Justificación} \\ \hline
\endfirsthead
\hline
\textbf{Entidad} & \textbf{Llaves candidatas} & \textbf{Forma} & \textbf{Justificación} \\ \hline
\endhead
\textit{Ranura} & (\texttt{id\_\allowbreak{}ranura}); (\texttt{codigo}) & FNBC & Llave simple; \texttt{codigo} depende solo de ella y es llave candidata. \\ \hline
\textit{Objeto\-Cosmetico} & (\texttt{id\_\allowbreak{}objeto}); (\texttt{codigo}) & FNBC & Llave simple; \texttt{codigo} es llave candidata. De la ranura se guarda solo su llave, no sus datos. (\texttt{id\_\allowbreak{}objeto}, \texttt{id\_\allowbreak{}ranura}) es superllave, no llave candidata: existe para las llaves foráneas compuestas. \\ \hline
\textit{Division} & (\texttt{id\_\allowbreak{}division}); (\texttt{codigo}); (\texttt{elo\_\allowbreak{}minimo}) & FNBC & Llave simple; \texttt{codigo} y \texttt{elo\_\allowbreak{}minimo} son llaves candidatas. El orden de las divisiones es el de \texttt{elo\_\allowbreak{}minimo}; guardar además una posición repetiría ese dato. \\ \hline
\textit{Nivel\-IA} & (\texttt{id\_\allowbreak{}nivel\_\allowbreak{}ia}); (\texttt{codigo}) & FNBC & Llave simple; \texttt{codigo} es llave candidata. \\ \hline
\textit{Leccion} & (\texttt{id\_\allowbreak{}leccion}); (\texttt{codigo}); (\texttt{orden}) & FNBC & Llave simple; \texttt{codigo} y \texttt{orden} son llaves candidatas. \\ \hline
\textit{Tipo\-Caja} & (\texttt{id\_\allowbreak{}tipo\_\allowbreak{}caja}); (\texttt{codigo}) & FNBC & Llave simple; \texttt{codigo} es llave candidata. \\ \hline
\textit{Tipo\-Caja\-Objeto} & (\texttt{id\_\allowbreak{}tipo\_\allowbreak{}caja}, \texttt{id\_\allowbreak{}objeto}) & FNBC & Llave compuesta; \texttt{probabilidad} depende del par completo: el mismo objeto tiene probabilidades distintas en cajas distintas. \\ \hline
\textit{Historial\-Nickname} & (\texttt{id\_\allowbreak{}usuario}) & FNBC & La llave es la llave foránea: relación 1:0..1 con \textit{Usuario}. El nombre vigente se lee de \textit{Usuario}; aquí queda solo el que ya no está (\mbox{CU-13}). No guarda el nombre nuevo: como el cambio es único, sería siempre una copia del \texttt{nickname} de \textit{Usuario}. \\ \hline
\textit{Avatar} & (\texttt{id\_\allowbreak{}avatar}) & FNBC & Llave simple. \\ \hline
\textit{Enlace\-Red} & (\texttt{id\_\allowbreak{}enlace}); (\texttt{id\_\allowbreak{}usuario}, \texttt{orden}) & FNBC & Llave sustituta; (\texttt{id\_\allowbreak{}usuario}, \texttt{orden}) es llave candidata. \\ \hline
\textit{Usuario\-Objeto} & (\texttt{id\_\allowbreak{}usuario}, \texttt{id\_\allowbreak{}objeto}) & FNBC & Llave compuesta; \texttt{fecha\_\allowbreak{}obtencion} y \texttt{origen} dependen del par usuario-objeto completo. Es la relación Posee del diagrama. \\ \hline
\textit{Equipamiento} & (\texttt{id\_\allowbreak{}usuario}, \texttt{id\_\allowbreak{}ranura}); (\texttt{id\_\allowbreak{}usuario}, \texttt{id\_\allowbreak{}objeto}) & 3FN & Llaves candidatas (\texttt{id\_\allowbreak{}usuario}, \texttt{id\_\allowbreak{}ranura}) e (\texttt{id\_\allowbreak{}usuario}, \texttt{id\_\allowbreak{}objeto}). La dependencia \texttt{id\_\allowbreak{}objeto} $\rightarrow$ \texttt{id\_\allowbreak{}ranura} parte de algo que no es superllave, pero \texttt{id\_\allowbreak{}ranura} es atributo primo: cumple 3FN y no FNBC. La llave foránea compuesta impide que ranura y objeto se contradigan. \\ \hline
\textit{Participacion\-Objeto} & (\texttt{id\_\allowbreak{}partida}, \texttt{id\_\allowbreak{}usuario}, \texttt{id\_\allowbreak{}ranura}); (\texttt{id\_\allowbreak{}partida}, \texttt{id\_\allowbreak{}usuario}, \texttt{id\_\allowbreak{}objeto}) & 3FN & Mismo caso que \textit{Equipamiento}: \texttt{id\_\allowbreak{}objeto} $\rightarrow$ \texttt{id\_\allowbreak{}ranura} con \texttt{id\_\allowbreak{}ranura} dentro de la llave; cumple 3FN y no FNBC. \\ \hline
\textit{Enlace\-Espectador} & (\texttt{id\_\allowbreak{}enlace}); (\texttt{token\_\allowbreak{}hash}) & FNBC & Llave simple; \texttt{token\_\allowbreak{}hash} es llave candidata. Sin indicador de activo: el enlace vale mientras su \textit{Partida} está \texttt{EN\_\allowbreak{}CURSO} (\mbox{CU-34} \mbox{RN-06}), y un indicador repetiría ese estado. \\ \hline
\textit{Invitacion} & (\texttt{id\_\allowbreak{}invitacion}) & FNBC & Llave simple. \\ \hline
\textit{Historial\-Division} & (\texttt{id\_\allowbreak{}historial\_\allowbreak{}division}) & FNBC & Llave simple. \\ \hline
\textit{Caja} & (\texttt{id\_\allowbreak{}caja}) & FNBC & Llave simple. \\ \hline
\textit{Caja\-Contenido} & (\texttt{id\_\allowbreak{}caja}, \texttt{numero}) & FNBC & Llave compuesta. \texttt{convertido} pasó a calculada porque dependía de \texttt{monedas\_\allowbreak{}conversion}. \\ \hline
\textit{Movimiento\-Moneda} & (\texttt{id\_\allowbreak{}movimiento}) & FNBC & Llave simple; \texttt{tipo} discrimina qué origen aplica. \\ \hline
\textit{Amistad} & (\texttt{id\_\allowbreak{}usuario\_\allowbreak{}a}, \texttt{id\_\allowbreak{}usuario\_\allowbreak{}b}) & FNBC & Llave compuesta de par ordenado: el CHECK impide guardar la misma amistad en los dos sentidos. \\ \hline
\textit{Solicitud} & (\texttt{id\_\allowbreak{}solicitud}); (\texttt{id\_\allowbreak{}solicitante}, \texttt{id\_\allowbreak{}destinatario}) & FNBC & Llave sustituta; (\texttt{id\_\allowbreak{}solicitante}, \texttt{id\_\allowbreak{}destinatario}) es llave candidata. \\ \hline
\textit{Progreso\-Tutorial} & (\texttt{id\_\allowbreak{}usuario}, \texttt{id\_\allowbreak{}leccion}) & FNBC & Llave compuesta usuario-lección. \texttt{completada} pasó a calculada porque dependía de \texttt{fecha\_\allowbreak{}completada}. \\ \hline
\end{longtable}
\end{center}
